\chapter{実験2}

\section{本章の概要}

本章では、工程文脈依存性を評価する実験2と、実機への接続可能性を確認する実験3について述べる。いずれも実験設計と評価方針を中心に整理する。

\section{実験2の設定}

実験2の目的は、同じ画像・同じ不確実性に対して、工程文脈が変化した場合に、VLM が処理分岐を適切に切り替えられるかを評価することである。比較対象は、現在工程なし、履歴なし、次工程候補なしのアブレーション条件である。評価では、Branch Contrast Accuracy を中心に、工程文脈情報の有無が分岐切り替え性能へ与える影響を確認する。

\section{評価の観点}

実験完了後には、提案手法が同一画像に対して工程文脈だけを変えた場合でも、文脈に応じて分岐を切り替えられるかを確認する予定である。特に、履歴なし条件では \texttt{act} と \texttt{reobserve} の区別が難しくなり、現在工程なし条件では \texttt{defer} の判断が不安定になる可能性があるため、これらの誤分類傾向を重点的に確認する。

\section{実験3：実機分岐デモ}

実験3の目的は、提案手法による診断結果が、実機ロボットの安全な分岐判断に接続できるかを確認することである。具体的には、\texttt{act} すべき場面で予定 primitive へ進めるか、\texttt{reobserve} すべき場面で追加観測へ移れるか、\texttt{ask\_user} すべき場面で人間確認へ回せるか、\texttt{stop} すべき場面で誤実行を避けられるか、\texttt{defer} すべき場面で過剰確認を避けられるかを評価する。

\begin{table}[tbp]
\centering
\caption{実機デモにおける分岐別成功基準}
\label{tab:demo_success}
\begin{tabularx}{\linewidth}{lX}
\toprule
分岐 & 成功基準 \\
\midrule
act & 予定された primitive が安全ゲートを通過し、動作ログが保存される \\
reobserve & 追加観測が実行され、再診断用の画像とログが保存される \\
ask\_user & 確認要求文が生成され、回答前は実行停止状態を維持する \\
stop & 危険または不明条件で動作を抑制し、停止ログを保存する \\
defer & 現在工程を継続し、不確実性を後工程の確認対象として記録する \\
\bottomrule
\end{tabularx}
\end{table}

評価では、VLM 単独の直接判断と比較し、False Act Rate と unsafe non-stop rate が低下するかを確認する。

\section{考察}

実験2で高い Branch Contrast Accuracy が得られる傾向が確認されれば、長期タスク中の不確実性診断には、画像理解だけでなく、作業履歴と現在工程の利用が重要であると解釈できる。さらに、実験3で false act rate や unsafe non-stop rate の低下が確認できれば、提案手法が Physical AI における行動前確認の仕組みとして有効である可能性を示せる。一方で、\texttt{ask\_user} や \texttt{stop} を過剰に選択する傾向が確認された場合は、実用上のテンポやユーザ負担とのトレードオフが課題となる。
